% !TeX root = part9.tex
\documentclass[../final.tex]{subfiles}
\begin{document}
% Конспект: Desktop/11.jpg, 22.jpg, 33.jpg; дата в тетради 30.10.
% Во второй строке системы для u, v в конспекте стоит 1+cos(ka) ---
% описка, исправлено на 1-cos(ka), итоговый закон дисперсии тот же.
% Интегралы для газа магнонов, записанные в конце, доведены до ответа.
\practice{9}{30.10.2026}{Спиновые волны}

\rtheory{Ферромагнитная цепочка}
Полагаем \rassumption{$k_B=1$}. Рассмотрим цепочку из $N$ спинов
$\boldsymbol S_p$ с периодом $a$ и периодическими граничными
условиями $\boldsymbol S_{N+1}=\boldsymbol S_1$. Обменное
взаимодействие соседей
\[
 U=-2J\sum_{p=1}^N\boldsymbol S_p\cdot\boldsymbol S_{p+1},
 \qquad J>0.
\]
При $T=0$ все спины параллельны оси $z$, и энергия основного состояния
\[
 U_0=-2JNS^2.
\]
При \rassumption{низкой температуре} спины лишь слегка отклоняются
от оси:
\[
 S_p^z\simeq S,\qquad |S_p^x|,\,|S_p^y|\ll S.
\]
Спины рассматриваем как классические векторы длины $S$.

\subsection*{Уравнение движения спина}
Магнитный момент атома $\boldsymbol\mu_p=-g\mu_B\boldsymbol S_p$.
Энергию связей спина $p$ с обоими соседями запишем как энергию
момента в эффективном поле:
\[
 U_p=-2J\boldsymbol S_p\cdot
 \left(\boldsymbol S_{p-1}+\boldsymbol S_{p+1}\right)
 =-\boldsymbol\mu_p\cdot\boldsymbol B_p,\qquad
 \boldsymbol B_p=-\frac{2J}{g\mu_B}
 \left(\boldsymbol S_{p-1}+\boldsymbol S_{p+1}\right).
\]
Скорость изменения момента импульса $\hbar\boldsymbol S_p$ равна
моменту сил $\boldsymbol\mu_p\times\boldsymbol B_p$:
\[
 \hbar\frac{d\boldsymbol S_p}{dt}
 =2J\,\boldsymbol S_p\times
 \left(\boldsymbol S_{p-1}+\boldsymbol S_{p+1}\right).
\]
Раскрываем векторное произведение и оставляем члены первого порядка
по малым поперечным компонентам:
\begin{align*}
 \hbar\frac{dS_p^x}{dt}&=2JS\left(2S_p^y-S_{p-1}^y-S_{p+1}^y\right),\\
 \hbar\frac{dS_p^y}{dt}&=-2JS\left(2S_p^x-S_{p-1}^x-S_{p+1}^x\right),\\
 \hbar\frac{dS_p^z}{dt}&=0.
\end{align*}

\subsection*{Закон дисперсии}
Ищем решение в виде бегущей волны:
\[
 S_p^x=\operatorname{Re}\left(u\,e^{i(kpa-\omega t)}\right),\qquad
 S_p^y=\operatorname{Re}\left(v\,e^{i(kpa-\omega t)}\right),\qquad
 S_p^z=S.
\]
Подставляя и учитывая $2-e^{ika}-e^{-ika}=2(1-\cos ka)$, получаем
\[
 \begin{cases}
  -i\omega\hbar\,u=4JS(1-\cos ka)\,v,\\
  -i\omega\hbar\,v=-4JS(1-\cos ka)\,u.
 \end{cases}
\]
Нетривиальное решение существует, когда определитель системы равен нулю:
\[
 \begin{vmatrix}
  i\omega\hbar & 4JS(1-\cos ka)\\
  -4JS(1-\cos ka) & i\omega\hbar
 \end{vmatrix}=0
 \quad\Longrightarrow\quad
 \hbar\omega=4JS(1-\cos ka)=8JS\sin^2\frac{ka}2.
\]
При этом $v=-iu$, то есть
\[
 S_p^x=u\cos(kpa-\omega t),\qquad S_p^y=u\sin(kpa-\omega t):
\]
каждый спин прецессирует вокруг оси $z$ по конусу, а фаза прецессии
соседних спинов сдвинута на $ka$. Эта волна называется
\rconcept{спиновой волной}. Из периодичности цепочки
$k=2\pi m/(Na)$, и в зоне Бриллюэна $-\pi/a<k\le\pi/a$ помещается
ровно $N$ волн.

В длинноволновом пределе $\lambda=2\pi/k\gg a$, то есть $ka\ll1$,
раскладываем косинус в ряд:
\[
 \hbar\omega\simeq2JSa^2k^2.
\]
В отличие от акустических фононов, частота квадратична по $k$.

\subsection*{Квантование амплитуды}
\begin{rbox}[left=16pt,right=16pt,top=12pt,bottom=10pt]{[]}
\centering
\begin{tikzpicture}[>=Latex,scale=1.15,every node/.style={text=rtext}]
    \coordinate (O) at (0,0);
    \coordinate (C) at (0,4);
    \coordinate (A) at ({2.2*cos(235)},{4+0.55*sin(235)});
    \coordinate (B) at ({2.2*cos(305)},{4+0.55*sin(305)});
    \coordinate (H) at (0,{4+0.55*sin(235)});
    \draw[densely dotted] (O)--(C);
    \draw[->] (C)--(0,5.1) node[above] {$z$};
    \draw[rc3,dashed] (C) ellipse [x radius=2.2,y radius=0.55];
    \draw[rc3,densely dotted] (O)--(-2.2,4) (O)--(2.2,4);
    \draw[rnotecolor] (C)--(A) node[midway,above left=-1pt] {$u$};
    \draw[rnotecolor] (C)--(B);
    \draw[rnotecolor] (A)--(B);
    \draw[->,thick,rconceptcolor] (O)--(A)
        node[pos=0.55,left] {$\boldsymbol S_p$};
    \draw[->,thick,rconceptcolor] (O)--(B)
        node[pos=0.55,right] {$\boldsymbol S_{p+1}$};
    \draw ({1.25*cos(109.58)},{1.25*sin(109.58)})
        arc[start angle=109.58,end angle=70.42,radius=1.25];
    \node at (0.2,1.5) {$\varphi$};
    \node[above right=-1pt] at (0.04,4.02) {$ka$};
    \fill (O) circle (1.5pt) node[below left] {$O$};
    \fill (A) circle (1.3pt) node[below left] {$A$};
    \fill (B) circle (1.3pt) node[below right] {$B$};
    \fill (H) circle (1.1pt) node[below] {$H$};
    \fill (C) circle (1.1pt) node[above left] {$C$};
\end{tikzpicture}
\captionof{figure}{Прецессия соседних спинов. Концы спинов движутся
по окружности радиуса $u$ со сдвигом фаз $ka$;
$\varphi$ --- угол между соседними спинами.}
\end{rbox}

Из треугольника $ACB$ с углом $ka$ при вершине $C$:
\[
 AH=u\sin\frac{ka}2,\qquad AB=2u\sin\frac{ka}2\simeq kau.
\]
Угол между соседними спинами находим из треугольника $AOH$:
\[
 \sin\frac\varphi2=\frac{AH}{AO}\simeq\frac{kau}{2S}
 \quad\Longrightarrow\quad
 \varphi\simeq\frac{kau}S.
\]
Поскольку $\boldsymbol S_p\cdot\boldsymbol S_{p+1}=S^2\cos\varphi$
для всех $p$,
\[
 U=-2JNS^2\cos\varphi
 \simeq-2JNS^2\left(1-\frac{\varphi^2}2\right)
 =U_0+JNk^2a^2u^2.
\]
Проекция каждого спина на ось уменьшается:
\[
 S_z=\sqrt{S^2-u^2}\simeq S-\frac{u^2}{2S}.
\]
\rassumption{Амплитуда $u$ мала}, так как мала температура.
В квантовой механике полная проекция $\sum_pS_p^z$ может принимать
только значения $NS-n$ с целым $n$. Поэтому
\[
 N(S-S_z)=\frac{Nu^2}{2S}=n
 \quad\Longrightarrow\quad
 u^2=\frac{2Sn}N.
\]
Подставляя в энергию, получаем
\[
 U=U_0+JNk^2a^2\frac{2Sn}N=U_0+n\cdot2JSa^2k^2=U_0+n\hbar\omega(k).
\]
\rresult{Энергия спиновой волны квантуется порциями $\hbar\omega(k)$}.
Квант спиновой волны называется \rconcept{магноном}; он несёт
квазиимпульс $\hbar k$ и уменьшает полную проекцию спина на единицу.

\rtheory{Газ магнонов}
Число магнонов $n$ в одной волне может быть любым, поэтому магноны ---
бозоны. Их число не сохраняется, и в равновесии
\[
 \mu=0,\qquad \overline n(\varepsilon)=\frac1{e^{\varepsilon/T}-1}.
\]
При $T\ll JS$ возбуждены только длинноволновые магноны с квадратичным
законом дисперсии:
\[
 \boldsymbol p=\hbar\boldsymbol k,\qquad
 \varepsilon=\frac{p^2}{2m_{\mathrm{eff}}},\qquad
 m_{\mathrm{eff}}=\frac{\hbar^2}{4JSa^2}.
\]
Для простой кубической решётки
$\hbar\omega=2JS\sum_{\boldsymbol\delta}(1-\cos\boldsymbol k\cdot\boldsymbol\delta)
\simeq2JSa^2k^2$ --- то же выражение, что и в цепочке, поэтому
дальше рассматриваем трёхмерный кристалл объёма $V$ из $N=V/a^3$ атомов.
У магнона одна поляризация, и
\[
 E=\int_0^\infty\frac{\varepsilon}{e^{\varepsilon/T}-1}
 \frac{V\,d^3p}{(2\pi\hbar)^3},\qquad
 N_m=\int_0^\infty\frac1{e^{\varepsilon/T}-1}
 \frac{V\,d^3p}{(2\pi\hbar)^3}.
\]
Верхний предел заменён бесконечностью: при $T\ll JS$ вклад
коротких волн экспоненциально мал.

\subsection*{Число магнонов и намагниченность}
Подставляя $x=\varepsilon/T$ и используя
$\int_0^\infty x^{s-1}/(e^x-1)\,dx=\Gamma(s)\zeta(s)$, находим
\[
 N_m=\zeta\left(\tfrac32\right)V
 \left(\frac{m_{\mathrm{eff}}T}{2\pi\hbar^2}\right)^{3/2}
 =\zeta\left(\tfrac32\right)N\left(\frac{T}{8\pi JS}\right)^{3/2}.
\]
Каждый магнон уменьшает полную проекцию спина на единицу, поэтому
намагниченность $M\propto NS-N_m$:
\[
 \frac{M(0)-M(T)}{M(0)}=\frac{N_m}{NS}
 =\frac{\zeta(3/2)}S\left(\frac{T}{8\pi JS}\right)^{3/2}
 \simeq\frac{0{,}0587}S\left(\frac{T}{2JS}\right)^{3/2}.
\]
Это \rconcept{закон Блоха}: \rresult{намагниченность убывает
как $T^{3/2}$}.

\subsection*{Энергия и теплоёмкость}
Аналогично
\[
 E=\frac32\zeta\left(\tfrac52\right)T\,V
 \left(\frac{m_{\mathrm{eff}}T}{2\pi\hbar^2}\right)^{3/2}
 =\frac32\zeta\left(\tfrac52\right)NT
 \left(\frac{T}{8\pi JS}\right)^{3/2}\propto T^{5/2},
\]
\[
 C_V=\frac{dE}{dT}
 =\frac{15}4\zeta\left(\tfrac52\right)N
 \left(\frac{T}{8\pi JS}\right)^{3/2}\propto T^{3/2}.
\]
\rresult{Магнонная теплоёмкость пропорциональна $T^{3/2}$}; при низких
температурах она преобладает над фононной ($\propto T^3$).
В обычных единицах $T$ заменяется на $k_BT$, а теплоёмкость умножается
на $k_B$.

\remark[Границы модели.]
Формулы справедливы при \rassumption{$T\ll JS$}, когда $N_m\ll NS$
и взаимодействием магнонов можно пренебречь. В одномерной и двумерной
моделях с $\varepsilon\propto p^2$ интеграл для $N_m$ расходится
на малых импульсах: $\int dp/p^2$ и $\int dp/p$. Поэтому в изотропной
цепочке и плёнке при $T>0$ дальний ферромагнитный порядок
\rnote{разрушается тепловыми спиновыми волнами}.
\end{document}
